% Part: many-valued-logic
% Chapter: three-valued-logics
% Section: lukasiewicz

\documentclass[../../../include/open-logic-section]{subfiles}

\begin{document}

\olfileid[id]{mvl}{inf}{god}

\olsection{Logika-Logika G\"odel}

\begin{editorial}
  Ini adalah ``rintisan'' singkat untuk suatu bagian mengenai logika G\"odel
  bernilai tak hingga.
\end{editorial}

\begin{defn}\ollabel{def:goedel} Logika G\"odel bernilai tak
hingga~$\LogGod[\infty]$ didefinisikan dengan menggunakan matriks berikut:
\begin{enumerate}
  \item Bahasa proposisional baku~$\Lang L_0$ dengan
  $\lfalse$, $\lnot$, $\land$, $\lor$, $\lif$.
  \item Himpunan nilai kebenaran~$V_\infty$.
  \item $1$ merupakan satu-satunya nilai kebenaran yang ditetapkan, yakni
  $V^+ = \{1\}$.
  \item Fungsi kebenaran diberikan oleh fungsi-fungsi berikut:
  \begin{align*}
    \tf{\lfalse} & = 0\\
    \tf{\lnot}[\LogGod](x) & = \begin{cases}
      1 & \text{jika } x =0\\
      0 & \text{selain itu}
    \end{cases}\\
    \tf{\land}[\LogGod](x,y) & = \min(x,y)\\
    \tf{\lor}[\LogGod](x,y) & = \max(x,y)\\
    \tf{\lif}[\LogGod](x,y) & = \begin{cases}
      1 & \text{jika } x \le y\\
      y & \text{selain itu.}
    \end{cases}
    \end{align*}
\end{enumerate}
Logika G\"odel bernilai~$m$ didefinisikan dengan cara yang sama, kecuali
$V = V_m$.
\end{defn}

\begin{prop}
  Logika $\LogGod[3]$ yang didefinisikan oleh
  \olref[thr][god]{defn:goedel} sama dengan $\LogGod[3]$ yang didefinisikan
  oleh \olref{def:goedel}.
\end{prop}

\begin{proof}
  Hal ini dapat dilihat dengan membandingkan tabel kebenaran untuk
  konektif-konektif yang diberikan dalam \olref[thr][god]{defn:goedel}
  dengan tabel kebenaran yang ditentukan oleh persamaan-persamaan dalam
  \olref{def:goedel}:
  \begin{center}
    \begin{tabular}{c|c} 
      $\tf{\lnot}[\LogGod[3]]$ & \\ 
      \hline  
      $1$ & $0$ \\ 
      $1/2$ & $0$ \\
      $0$ & $1$ 
    \end{tabular}
    \quad
    \begin{tabular}{c|ccc} 
      $\tf{\land}[\LogGod]$ & $1$ & $1/2$ & $0$ \\ 
      \hline 
      $1$ & $1$ & $1/2$ & $0$ \\ 
      $1/2$ & $1/2$ & $1/2$ & $0$\\ 
      $0$ & $0$ & $0$ & $0$ 
    \end{tabular}
    \\[2ex]
    \begin{tabular}{c|ccc} 
      $\tf{\lor}[\LogGod]$ & $1$ & $1/2$ & $0$ \\ 
      \hline 
      $1$ & $1$ & $1$ & $1$ \\ 
      $1/2$ & $1$ & $1/2$ & $1/2$ \\
      $0$ & $1$ & $1/2$ & $0$ 
    \end{tabular}
    \quad
    \begin{tabular}{c|ccc} 
      $\tf{\lif}[\LogGod]$ & $1$ & $1/2$ & $0$ \\ 
      \hline 
      $1$ & $1$ & $1/2$ & $0$ \\ 
      $1/2$ & $1$ & $1$ & $0$  \\ 
      $0$ & $1$ & $1$ & $1$ 
    \end{tabular}
  \end{center} 
\end{proof}

\begin{prop}\ollabel{prop:god-infty-m}
  Jika $\Gamma \Entails[\LogGod[\infty]] !B$, maka $\Gamma
  \Entails[\LogGod[m]] !B$ untuk setiap~$m \ge 2$.
\end{prop}

\begin{proof}
  Latihan.
\end{proof}

\begin{prob}
  Buktikan \olref[mvl][inf][god]{prop:god-infty-m}.
\end{prob}

Faktanya, kebalikannya juga berlaku.

Sebagaimana dalam~$\LogGod[3]$, semua !!{formula}s yang valid secara
intuisionistik merupakan tautologi dari~$\LogGod[\infty]$, dan contoh-contoh
non-tautologi yang sama juga merupakan non-tautologi dari~$\LogGod[\infty]$:
\begin{align*}
  & p \lor \lnot p && (p \lif q) \lif (\lnot p \lor q) \\
  & \lnot\lnot p \lif p && \lnot(\lnot p \land \lnot q) \lif (p \lor q) \\
  & ((p \lif q) \lif p) \lif p && \lnot(p \lif q) \lif (p \land \lnot q)
\end{align*}
Contoh !!{formula} yang tidak valid secara intuisionistik tetapi tetap
merupakan tautologi dari~$\LogGod[3]$, $(p \lif q) \lor (q \lif p)$, juga
merupakan tautologi dalam~$\LogGod[\infty]$. Faktanya, $\LogGod[\infty]$
dapat dicirikan sebagai logika intuisionistik yang diperoleh dengan
menambahkan skema $(!A \lif !B) \lor (!B \lif !A)$. Hasil ini dibuktikan oleh
Michael Dummett; karena itu, $\LogGod[\infty]$ sering disebut logika
G\"odel--Dummett~$\Log{LC}$.

\begin{prob}
  Tunjukkan bahwa $(p \lif q) \lor (q \lif p)$ merupakan tautologi
  dari~$\LogGod[\infty]$.
\end{prob}

\begin{prob}
  Tunjukkan bahwa $(p \lif q) \lor (q \lif r) \lor (r \lif s)$, yang
  merupakan tautologi dari~$\LogGod[3]$, bukan tautologi
  dari~$\LogGod[\infty]$.
\end{prob}

\end{document}
